% !TeX program = xelatex
% 华为杯研赛论文 LaTeX 模板（第二十三届 2026 格式规范）
%
% 编译顺序（必须用 XeLaTeX，不能用 pdflatex）：
%   xelatex main.tex
%   bibtex  main            % 仅在改用 BibTeX 文献时
%   xelatex main.tex
%   xelatex main.tex
%
% 官方硬性要求（见 references/paper-formatting.md）：
%   题目三号黑体居中；一级标题四号黑体居中；其余汉字小四宋体；单倍行距；
%   页码从摘要页起位于页脚中部；不得有页眉；不得出现身份信息。

\documentclass[UTF8,a4paper,zihao=-4]{ctexart}

% ---------- 页面：取自官方附件3 论文模板 ----------
% 模板实测：A4 210×297 mm；页边距 上30.0 下17.5 左右22.5 mm；页眉距15.0、页脚距17.5 mm
\usepackage{geometry}
\geometry{a4paper, top=3.0cm, bottom=1.75cm, left=2.25cm, right=2.25cm, footskip=0.9cm}

% ---------- 西文与数字：Times New Roman ----------
\usepackage{fontspec}
\setmainfont{Times New Roman}

% ---------- 行距：官方要求单倍 ----------
\usepackage{setspace}
\singlespacing

% ---------- 标题格式：一级标题四号黑体居中 ----------
\usepackage{titlesec}
\titleformat{\section}{\centering\heiti\zihao{4}}{\thesection}{0.5em}{}
\titleformat{\subsection}{\heiti\zihao{-4}}{\thesubsection}{0.5em}{}
\titleformat{\subsubsection}{\heiti\zihao{-4}}{\thesubsubsection}{0.5em}{}
\titlespacing*{\section}{0pt}{1.4ex plus .3ex}{0.9ex}
\titlespacing*{\subsection}{0pt}{1.1ex plus .2ex}{0.6ex}
\setcounter{secnumdepth}{3}

% ---------- 图表与公式编号：按章编号，图题在下、表题在上 ----------
\usepackage{caption}
\usepackage{booktabs}
\usepackage{graphicx}
\usepackage{amsmath,amssymb}
\usepackage{array}
\captionsetup{font=small, labelsep=quad, justification=centering}
\captionsetup[figure]{position=below, skip=6pt}
\captionsetup[table]{position=above, skip=6pt}
\numberwithin{figure}{section}
\numberwithin{table}{section}
\numberwithin{equation}{section}

% ---------- 页眉页脚：无页眉，页脚居中页码 ----------
\usepackage{fancyhdr}
\pagestyle{fancy}
\fancyhf{}
\fancyfoot[C]{\zihao{5}\thepage}
\renewcommand{\headrulewidth}{0pt}
\renewcommand{\footrulewidth}{0pt}

% ---------- 附录代码 ----------
\usepackage{listings}
\usepackage{xcolor}
\lstset{
  basicstyle=\ttfamily\zihao{5},
  breaklines=true,
  columns=flexible,
  showstringspaces=false,
  frame=none,
  numbers=left,
  numberstyle=\tiny\color{gray},
  keywordstyle=\color{black},
  commentstyle=\color{gray},
}

% ---------- 超链接（不影响打印） ----------
\usepackage[hidelinks]{hyperref}

% ============================================================
\begin{document}

% ---------- 摘要页：题目、摘要、关键词同页 ----------
\begin{center}
  {\heiti\zihao{3} 基于 XX 模型的 XX 问题研究}   % 题目：三号黑体居中
\end{center}
\vspace{1em}

\begin{center}
  {\heiti\zihao{4} 摘\quad 要}
\end{center}

% 摘要正文：小四宋体；含建模思路、主要方法、模型、结果与结论、创新点
摘要正文。针对问题一，……。针对问题二，……。针对问题三，……。
本文的创新点在于……。结论表明……。

\vspace{0.6em}
\noindent{\heiti 关键词：}{\songti 关键词一；关键词二；关键词三；关键词四}

\newpage

% ---------- 正文 ----------
\section{问题重述}

\subsection{问题背景}

正文使用小四号宋体、单倍行距。所有汉字均为小四号宋体。

\section{模型假设与符号说明}

符号表建议使用三线表，见表~\ref{tab:symbols}。

\begin{table}[htbp]
  \centering
  \caption{符号说明}
  \label{tab:symbols}
  \begin{tabular}{cll}
    \toprule
    符号 & 含义 & 单位 \\
    \midrule
    $x$        & 决策变量 & --- \\
    $\rho$     & 密度     & kg/m$^3$ \\
    \bottomrule
  \end{tabular}
\end{table}

\section{问题一：模型建立与求解}

行内公式写作 $E = mc^2$；独立公式按章编号，如式~\eqref{eq:sample}。

\begin{equation}
  \label{eq:sample}
  \min_{\boldsymbol{x}} \; f(\boldsymbol{x}) \quad \text{s.t.} \quad
  \boldsymbol{A}\boldsymbol{x} \le \boldsymbol{b}, \;
  \boldsymbol{x} \ge \boldsymbol{0}
\end{equation}

式中 $f(\cdot)$ 为目标函数，$\boldsymbol{A}$ 为约束矩阵。

图表必须在正文被引用并解释，例如图~\ref{fig:demo} 所示。

\begin{figure}[htbp]
  \centering
  \includegraphics[width=0.8\textwidth]{figures/demo.pdf}
  \caption{示例结果图}
  \label{fig:demo}
\end{figure}

\section{模型评价与推广}

\section*{参考文献}
\addcontentsline{toc}{section}{参考文献}

% 官方三种表述方式：书籍 / 期刊 / 网上资源。
% 正文引用用方括号，如 [1][3]；引用书籍必须指出页码。
\begin{thebibliography}{99}
  \bibitem{ref1} 姜启源，数学模型，北京：高等教育出版社，12-45，2011。
  \bibitem{ref2} 王某，李某，基于改进遗传算法的路径规划，系统工程理论与实践，42(3)：512-521，2022。
  \bibitem{ref3} 中国研究生数学建模竞赛组委会，竞赛章程，https://cpipc.acge.org.cn，2026年9月16日。
\end{thebibliography}

\appendix
\section*{附录}

源代码按题目要求另行提交竞赛系统；此处仅放必要的代码索引与说明。

\end{document}
